\documentclass[12pt]{article}
\usepackage[utf8]{inputenc}
\usepackage[T1]{fontenc}
\usepackage[margin=1in]{geometry}
\usepackage{setspace}
\usepackage{xcolor}
\usepackage[colorlinks=true, urlcolor=blue!60!black, linkcolor=blue!60!black]{hyperref}
\setlength{\parindent}{0em}
\setlength{\parskip}{0.6em}
\renewcommand{\baselinestretch}{1.0}
\begin{document}
\begin{center}
{\LARGE \textit{Prigg v. Pennsylvania}, 41 U.S. 539 (1842)\par}
\vspace{0.5em}
{\large Case Brief --- Two-Page Edition\par}
\end{center}

\section{Memory Jogger}
Pennsylvania's anti-kidnapping law fell: the Fugitive Slave Clause gives a self-executing right of recapture, and Congress's power over rendition is exclusive (\href{https://supreme.justia.com/cases/federal/us/41/539/}{Justia}).

\section{Detailed Case Facts}
Margaret Morgan, held as a slave in Maryland, moved to Pennsylvania in 1832. In 1837 Edward Prigg, an agent for her owner, forcibly carried Morgan and her children to Maryland after a local magistrate declined to act. Pennsylvania's 1826 statute made such forcible removal a felony.

\section{Procedural History}
In an agreed test case, Prigg was convicted of kidnapping on a special verdict, and the Pennsylvania Supreme Court affirmed \textit{pro forma}.

\section{Judicial Votes}
Reversed 8--1 with seven opinions. Story wrote for the Court; Taney, Thompson, Baldwin, and Daniel concurred in the judgment on differing grounds; McLean dissented.

\section{Holding}
A slaveholder may peaceably recapture a fugitive in any State, and state laws obstructing recaption are void. Congress's power over rendition is exclusive, but States cannot be compelled to enforce the federal Act.

\section{Analysis of Opinions}
\textbf{Story:} the Clause created a ``positive, unqualified right''; Congress's implied power to enforce it is exclusive, barring even ``auxiliary'' state laws; state magistrates may act only ``unless prohibited by state legislation.'' \textbf{Taney:} the Clause bars only state laws that impair the right, not those that support it. \textbf{McLean:} a free State presumes everyone free, so recaption required a federal certificate.

\section{Examples: Future Applications}
\textbf{Same-Side:} a State's extra certificates for federally regulated railroads are preempted. \textbf{Opposite-Side:} a State may withdraw its officers' help from a federal program. \textbf{Fence-Sitter:} state penalties added ``to aid'' federal hazardous-materials rules.

\section{Critique}
Story inferred both an unenumerated power and its exclusivity from textual silence. Finkelman rejects the ``triumph of freedom'' defense, showing that Story stripped free Black Northerners of protection against kidnapping.

\section{Key Quotations}
\begin{quote}
``Upon this ground, we have not the slightest hesitation in holding that, under and in virtue of the Constitution, the owner of a slave is clothed with entire authority, in every State in the Union, to seize and recapture his slave whenever he can do it without any breach of the peace or any illegal violence. In this sense and to this extent, this clause of the Constitution may properly be said to execute itself, and to require no aid from legislation, state or national.''
\end{quote}
\noindent --- Justice Story, for the Court, 41 U.S. at 613.

\end{document}
